\chapter{{\fontsize{14pt}{21pt}\selectfont\MakeUppercase{Проблема и целевая аудитория}}}
\label{ch:problem}

\section{Какую проблему решает проект}

В~условиях цифровой трансформации малый и средний бизнес (далее~--- МСБ) сталкивается с~необходимостью поддержания актуального и согласованного цифрового присутствия на множественных платформах: картографических сервисах (Яндекс.Карты, 2ГИС), социальных сетях (VK), мессенджерах (Telegram), агрегаторах товаров и услуг (Avito). Ручное управление этими каналами, включая публикацию контента, обновление информации и мониторинг обратной связи, требует значительных временных и трудозатрат.

Это приводит к~несвоевременному обновлению данных, противоречивой информации на разных площадках и, как следствие, к~ухудшению клиентского опыта, потере потенциальных клиентов и репутационным рискам. Дополнительной сложностью является гетерогенность интеграционных интерфейсов: одни платформы предоставляют официальные API (VK, Telegram), тогда как другие, включая лидирующие на российском рынке геосервисы (Яндекс.Бизнес, 2ГИС), не имеют публичных API для программного управления карточкой организации. Существующие комплексные SaaS-решения зачастую являются избыточными и финансово недоступными для малого бизнеса.

Задача автоматизации управления цифровым присутствием МСБ на нескольких внешних площадках включает поддержание актуальности данных и обработку обратной связи. Проблема формализуется так: данные о~компании и коммуникации с~клиентами распределены по нескольким платформам; обновления выполняются вручную или разрозненными инструментами, что повышает риск ошибок и несогласованности; отсутствие единого контура контроля и журналирования затрудняет выявление <<истины>> по бизнес-данным.

\section{У кого возникает проблема (целевая аудитория)}

Проблема возникает у субъектов малого и среднего бизнеса (МСБ), которые ведут деятельность с~офлайн-точкой присутствия и нуждаются в~актуальном отображении своей информации в~цифровых каналах: на картах, в~социальных сетях и мессенджерах.

В~качестве типового примера целевой аудитории рассматривается \textbf{кофейня}. По данным РОМИР, заведения общественного питания~--- наиболее частая категория поиска в~офлайне (84~\%), а~геосервисы используют для уточнения графика работы (48~\%)~\cite{src:35}. Для кофейни критичны корректные часы работы, адрес и маршрут, быстрые ответы на вопросы и отзывы. Кофейня представляет собой малое предприятие с~ограниченным штатом, для которого выделение отдельного сотрудника на управление цифровым присутствием экономически нецелесообразно; задача автоматизации особенно актуальна именно для данного сегмента.

Целевой рынок в~рамках проекта~--- прежде всего Россия (VK, Telegram, Яндекс.Бизнес, 2ГИС); архитектура является универсальной и допускает адаптацию к~другим странам путём замены набора агентов (например, Google Business Profile для международных рынков).
